% !PRINTED 184
\weberpage{184}
%<<CONT
\weq{1}{\begin{aligned}
\varphi(x) &= Y_1^2 + Y_2^2 + \cdots Y_\nu^2 - Y_1'^2 - Y_2'^2 - \cdots - Y_{\nu'}'^2\\
&= Z_1^2 + Z_2^2 + \cdots Z_\mu^2 - Z_1'^2 - Z_2'^2 - \cdots - Z_{\mu'}'^2
\end{aligned}}
zwei Zerlegungen von $\varphi(x)$ in Quadrate. Es seien also $Y_1, Y_2 \ldots Y_\nu, Y_1', Y_2' \ldots Y_{\nu'}'$ homogene lineare Functionen von $x$, zwischen denen keine lineare Relation mit constanten Co\"efficienten besteht und also $\nu + \nu' \leqq m$; dieselben Voraussetzungen werden über die Functionen $Z_1, Z_2 \ldots Z_\mu, Z_1', Z_2' \ldots Z_{\mu'}'$ gemacht, so dass auch $\mu + \mu' \leqq m$ ist.

Angenommen, es sei
\[
\nu < \mu,
\]
dann können wir die Variablen $x$ den linearen Gleichungen
\weq{2}{\begin{aligned}
Y_1 &= 0, \quad Y_2 = 0 \ldots Y_\nu = 0\\
Z_1' &= 0, \quad Z_2' = 0 \ldots Z_{\mu'}' = 0
\end{aligned}}
unterwerfen, deren Anzahl $\nu + \mu'$ kleiner als $m$ ist.

Wenn nun die sämmtlichen Functionen $Z_1, Z_2 \ldots Z_\mu$ linear abhängig wären von den Functionen $Y_1, Y_2 \ldots Y_\nu, Z_1', Z_2' \ldots Z_{\mu'}'$, so könnte man aus diesen $\mu$ Gleichungen durch Elimination der $\nu$ Variablen $Y_1, Y_2 \ldots Y_\nu$ eine lineare Gleichung zwischen den $Z_1, Z_2 \ldots Z_\mu, Z_1' \ldots Z_{\mu'}'$ herleiten, die nach Voraussetzung \gesperrt{nicht} bestehen soll. Es ist also das Verschwinden sämmtlicher $Z_1, Z_2 \ldots Z_\mu$ nicht eine nothwendige Folge der Gleichungen (2), und wir können zu den Gleichungen (2) noch eine nicht homogene hinzufügen, etwa
\[
Z_1 = 1,
\]
dann haben wir für die $m$ Unbekannten $x$ ein System von $m$ oder weniger Gleichungen, die von einander unabhängig sind, sich also nicht widersprechen.

Dann ergiebt aber die zweite Darstellung (1) für $\varphi(x)$ einen positiven Werth, während die erste einen negativen oder verschwindenden Werth giebt, worin ein Widerspruch liegt. Es folgt also
\[
\nu \geqq \mu,
\]
und da man ebenso schliesst
\[
\mu \geqq \nu,
\]
so bleibt nur $\nu = \mu$ übrig. In gleicher Weise kann man beweisen, dass $\nu' = \mu'$ ist, wodurch unser Satz vollständig bewiesen ist.
